% Vietnamese translation for OI Public Library
% Source-PDF: IOI中国国家候选队论文/2007/day2/9.周冬《生成树的计数及其应用》.pdf
\documentclass[11pt]{article}
\usepackage{oipl-vn}

\newcommand{\Oh}{\mathcal{O}}
\newcommand{\mat}[1]{\begin{pmatrix}#1\end{pmatrix}}
\newcommand{\minor}[3]{#1\!\left(\begin{matrix}#2\\#3\end{matrix}\right)}

\title{Đếm cây khung và ứng dụng}
\author{Zhou Dong\\An Huy}
\date{}

\begin{document}
\maketitle

\origtitle{Counting spanning trees and its applications}
\sourcepath{IOI China candidate team papers / 2007 / day 2 / Zhou Dong}

\begin{abstract}
Trong các kỳ thi tin học, ta thường gặp các bài toán tối ưu hóa liên quan đến
cây khung, chẳng hạn cây khung nhỏ nhất, nhưng ít khi gặp bài toán đếm cây
khung và các vấn đề liên quan. Thật ra, đếm cây khung là một chủ đề rất có ý
nghĩa và có ứng dụng rộng rãi trong nhiều phương diện. Bài viết bắt đầu từ
một ví dụ trong thi tin học: trước hết giới thiệu một thuật toán quy hoạch
động có độ phức tạp mũ, sau đó trình bày các khái niệm và tính chất cơ bản
của định thức. Trên cơ sở đó ta đưa vào định lý Matrix-Tree, đồng thời so
sánh với một bài toán toán học để làm rõ tư tưởng toán học ẩn trong định lý.
Cuối cùng, bài viết dùng một vài ví dụ để giới thiệu ứng dụng của việc đếm cây
khung trong thi tin học và đưa ra tổng kết.
\end{abstract}

\paragraph{Từ khóa.}
Đếm cây khung; định lý Matrix-Tree.

\section{Bài toán được đặt ra}

\subsection*{Ví dụ 1: Highways (SPOJ p104)}

Một quốc gia có \(n\) thành phố, đánh số từ \(1\) đến \(n\). Giữa một số cặp
thành phố có thể xây đường cao tốc. Bây giờ cần chọn xây một số đường cao tốc
để tạo thành một mạng giao thông. Nhiệm vụ là tính có bao nhiêu phương án sao
cho giữa hai thành phố bất kỳ có đúng một đường đi.

Quy mô dữ liệu: \(1\le n\le 12\).

\subsubsection*{Phân tích}

Ta có thể chuyển bài toán sang mô hình đồ thị. Vì giữa hai đỉnh bất kỳ có đúng
một đường đi, kết quả cuối cùng phải là một cây khung của đồ thị ban đầu. Do
đó bài toán trở thành: cho một đồ thị vô hướng \(G\), hãy tính số cây khung
\(t(G)\) của nó. Phải làm thế nào?

Qua phân tích, có thể thu được một thuật toán quy hoạch động với độ phức tạp
\(\Oh(3^n n^2)\). Với giới hạn nhỏ của bài gốc, cách này đủ dùng. Nhưng khi
\(n\) lớn hơn thì không còn khả thi. Có thuật toán tốt hơn không? Câu trả lời
là có. Trước khi giới thiệu thuật toán, ta cần học một số kiến thức chuẩn bị.

\section{Kiến thức chuẩn bị}

Dưới đây ta giới thiệu một công cụ đại số quan trọng: định thức. Để định nghĩa
định thức, trước hết ta xét khái niệm hoán vị.

\subsection{Hoán vị}

\paragraph{Định nghĩa 1.}
Một dãy có thứ tự \(i_1i_2\cdots i_n\) tạo bởi các số \(1,2,\ldots,n\) được gọi
là một hoán vị của \(1,2,\ldots,n\).

Theo định nghĩa, \(i_1,i_2,\ldots,i_n\) biểu diễn \(n\) số tự nhiên khác nhau,
và mỗi số trong đó đều thuộc tập
\[
  S_n=\{1,2,\ldots,n\}.
\]
Nói cách khác,
\[
  1\mapsto i_1,\quad 2\mapsto i_2,\quad \ldots,\quad n\mapsto i_n
\]
định nghĩa một song ánh từ \(S_n\) đến chính nó. Song ánh này có thể ký hiệu là
\[
  \mat{1 & 2 & \cdots & n\\ i_1 & i_2 & \cdots & i_n},
\]
và được gọi là một phép thế. Song ánh trên cũng có thể viết lại dưới dạng
\[
  j_1\mapsto i_{j_1},\quad j_2\mapsto i_{j_2},\quad \ldots,\quad
  j_n\mapsto i_{j_n},
\]
trong đó \(j_1j_2\cdots j_n\) là một hoán vị của \(1,2,\ldots,n\). Vì vậy song
ánh này cũng có thể ký hiệu là
\[
  \mat{j_1 & j_2 & \cdots & j_n\\
       i_{j_1} & i_{j_2} & \cdots & i_{j_n}}.
\]
Do đó, với mọi hoán vị \(j_1j_2\cdots j_n\) của \(1,2,\ldots,n\), ta định nghĩa
\[
  \mat{1 & 2 & \cdots & n\\ i_1 & i_2 & \cdots & i_n}
  =
  \mat{j_1 & j_2 & \cdots & j_n\\
       i_{j_1} & i_{j_2} & \cdots & i_{j_n}}.
\]

Lấy một hoán vị \(i_1i_2\cdots i_n\). Nếu đổi chỗ hai số tự nhiên kề nhau
\(i_{j-1}\) và \(i_j\), ta nhận được hoán vị mới
\[
  i_1i_2\cdots i_{j-2}i_ji_{j-1}i_{j+1}\cdots i_n.
\]
Hoán vị này gọi là hoán vị sau một phép đổi chỗ của hoán vị ban đầu, và thao
tác đó gọi là một phép đổi chỗ. Rõ ràng mọi hoán vị đều có thể biến thành hoán
vị chuẩn \(12\cdots n\) sau một số phép đổi chỗ. Dù chọn cách đổi chỗ nào, với
hoán vị \(i_1i_2\cdots i_n\) cho trước, số lần đổi chỗ có tính chất quan trọng
sau.

\paragraph{Định lý 1.}
Khi biến một hoán vị bất kỳ \(i_1i_2\cdots i_n\) thành hoán vị chuẩn
\(12\cdots n\) bằng các phép đổi chỗ, tính chẵn lẻ của số phép đổi chỗ cần dùng
không phụ thuộc vào cách đổi chỗ.

Dựa vào định lý này, ta có định nghĩa sau.

\paragraph{Định nghĩa 2.}
Một hoán vị \(i_1i_2\cdots i_n\) được gọi là hoán vị chẵn (tương ứng lẻ) nếu
tồn tại một cách dùng số chẵn (tương ứng lẻ) phép đổi chỗ để biến nó thành
hoán vị chuẩn \(12\cdots n\).

Giả sử hoán vị \(i_1i_2\cdots i_n\) biến thành \(12\cdots n\) sau \(t\) phép đổi
chỗ. Khi đó giá trị \((-1)^t\) không phụ thuộc vào cách đổi chỗ. Ta viết giá
trị này là
\[
  \delta\!\binom{12\cdots n}{i_1i_2\cdots i_n}=(-1)^t
  =
  \begin{cases}
    1, & \text{khi } i_1i_2\cdots i_n
          \text{ là hoán vị chẵn của } 1,2,\ldots,n,\\
   -1, & \text{khi } i_1i_2\cdots i_n
          \text{ là hoán vị lẻ của } 1,2,\ldots,n.
  \end{cases}
\]

Một hoán vị của \(n\) số tự nhiên cố định \(1,2,\ldots,n\) có thể xem là một
song ánh từ \(S_n\) đến chính nó. Ta mở rộng khái niệm này cho một tập bất kỳ
\[
  S=\{a_1,a_2,\ldots,a_n\}.
\]
Với song ánh từ \(S\) đến chính nó
\[
  a_1\mapsto a_{i_1},\quad a_2\mapsto a_{i_2},\quad \ldots,\quad
  a_n\mapsto a_{i_n},
\]
dãy \(a_{i_1}a_{i_2}\cdots a_{i_n}\) được gọi là một hoán vị của
\(a_1,a_2,\ldots,a_n\). Dễ thấy điều kiện cần và đủ để
\(a_{i_1}a_{i_2}\cdots a_{i_n}\) là một hoán vị của
\(a_1,a_2,\ldots,a_n\) là \(i_1i_2\cdots i_n\) là một hoán vị của
\(1,2,\ldots,n\).

Tương tự, tính chẵn lẻ của tổng số phép đổi chỗ cần dùng để biến
\(a_{i_1}a_{i_2}\cdots a_{i_n}\) thành hoán vị chuẩn
\(a_1a_2\cdots a_n\) cũng không phụ thuộc vào cách đổi chỗ. Do đó ta đặt
\[
  \delta\!\binom{a_1a_2\cdots a_n}{a_{i_1}a_{i_2}\cdots a_{i_n}}
  =(-1)^t
  =
  \begin{cases}
    1, & \text{khi } a_{i_1}a_{i_2}\cdots a_{i_n}
          \text{ là hoán vị chẵn của } a_1,a_2,\ldots,a_n,\\
   -1, & \text{khi } a_{i_1}a_{i_2}\cdots a_{i_n}
          \text{ là hoán vị lẻ của } a_1,a_2,\ldots,a_n.
  \end{cases}
\]
Ở đây \(t\) là tổng số phép đổi chỗ trong một cách biến
\(a_{i_1}a_{i_2}\cdots a_{i_n}\) thành \(a_1a_2\cdots a_n\).

\subsection{Định thức}

Định thức cấp một là hàm của một biến \(a_{11}\):
\[
  \det(a_{11})=a_{11}.
\]
Nó cũng có thể viết thành
\[
  \det(a_{11})=\sum_{\mat{1\\1}}
  \delta\!\binom{1}{1}a_{11}.
\]

Định thức cấp hai là hàm của bốn biến \(a_{11},a_{12},a_{21},a_{22}\):
\[
  \det\mat{a_{11} & a_{12}\\ a_{21} & a_{22}}
  =a_{11}a_{22}-a_{12}a_{21}.
\]
Cũng có thể viết nó thành
\[
  \det\mat{a_{11} & a_{12}\\ a_{21} & a_{22}}
  =
  \sum_{\mat{1&2\\ i_1&i_2}}
  \delta\!\binom{12}{i_1i_2}a_{1i_1}a_{2i_2}.
\]

Định thức cấp ba là hàm của chín biến \(a_{ij}\), \(i,j=1,2,3\):
\[
\begin{aligned}
  \det\mat{a_{11}&a_{12}&a_{13}\\
           a_{21}&a_{22}&a_{23}\\
           a_{31}&a_{32}&a_{33}}
  &=
  a_{11}a_{22}a_{33}+a_{12}a_{23}a_{31}+a_{13}a_{21}a_{32}\\
  &\quad
  -a_{13}a_{22}a_{31}-a_{12}a_{21}a_{33}-a_{11}a_{23}a_{32}.
\end{aligned}
\]
Tương tự, nó có thể viết là
\[
  \det\mat{a_{11}&a_{12}&a_{13}\\
           a_{21}&a_{22}&a_{23}\\
           a_{31}&a_{32}&a_{33}}
  =
  \sum_{\mat{1&2&3\\ i_1&i_2&i_3}}
  \delta\!\binom{123}{i_1i_2i_3}
  a_{1i_1}a_{2i_2}a_{3i_3}.
\]

Quan sát định nghĩa của định thức cấp một, hai, ba, ta nhận được định nghĩa
tổng quát của định thức cấp \(n\).

\paragraph{Định nghĩa 3.}
Định thức của ma trận cấp \(n\) \(A=(a_{ij})\) là một số thực, ký hiệu
\(\det A\), được định nghĩa bởi
\[
  \det A=
  \sum_{\mat{1&2&\cdots&n\\ i_1&i_2&\cdots&i_n}}
  \delta\!\binom{12\cdots n}{i_1i_2\cdots i_n}
  a_{1i_1}a_{2i_2}\cdots a_{ni_n}.
\]

Định thức thường được ký hiệu bằng các dạng tương đương sau:
\[
  \det A
  =
  \det\mat{a_{11}&\cdots&a_{1n}\\
           \vdots&       &\vdots\\
           a_{n1}&\cdots&a_{nn}}
  =
  \left|\begin{matrix}
    a_{11}&\cdots&a_{1n}\\
    \vdots&       &\vdots\\
    a_{n1}&\cdots&a_{nn}
  \end{matrix}\right|
  =|A|.
\]

Từ định nghĩa thấy rằng tính trực tiếp định thức bằng công thức trên là rất
khó; chỉ khi bậc thấp mới có thể làm vậy. Để tính toán hiệu quả, trước hết cần
suy ra một số tính chất cơ bản của định thức, rồi dùng các tính chất đó để
biến định thức phức tạp thành định thức đơn giản, hoặc biến định thức bậc cao
thành định thức bậc thấp.

\paragraph{Tính chất 1.}
Nếu \(A\) là ma trận cấp \(n\), thì
\[
  \det A^{\mathsf T}=\det A.
\]
Ý nghĩa quan trọng của tính chất này là: trong định thức, hàng và cột có vai
trò bình đẳng. Chính xác hơn, mọi tính chất đúng cho hàng cũng đúng cho cột,
và ngược lại.

\paragraph{Tính chất 2.}
Đổi chỗ hai hàng bất kỳ (hoặc hai cột bất kỳ) của định thức thì định thức đổi
dấu.

\paragraph{Hệ quả.}
Nếu định thức có hai hàng (hoặc hai cột) giống nhau thì giá trị của định thức
bằng \(0\).

\paragraph{Tính chất 3.}
Nhân một hàng (hoặc một cột) với số thực \(\lambda\), định thức nhận được bằng
\(\lambda\) lần định thức ban đầu.

\paragraph{Hệ quả.}
Nếu định thức có hai hàng (hoặc hai cột) tỉ lệ với nhau thì giá trị bằng
\(0\). Đặc biệt, nếu định thức có một hàng (hoặc một cột) toàn \(0\), giá trị
định thức bằng \(0\).

\paragraph{Tính chất 4.}
Định thức tuyến tính theo từng hàng. Chẳng hạn, nếu hàng thứ \(j\) là tổng của
hai hàng \((a_{j1},\ldots,a_{jn})\) và \((b_{j1},\ldots,b_{jn})\), thì
\[
\begin{aligned}
&\left|\begin{matrix}
a_{11}&\cdots&a_{1n}\\
\vdots&&\vdots\\
a_{j1}+b_{j1}&\cdots&a_{jn}+b_{jn}\\
\vdots&&\vdots\\
a_{n1}&\cdots&a_{nn}
\end{matrix}\right| \\
&\quad =
\left|\begin{matrix}
a_{11}&\cdots&a_{1n}\\
\vdots&&\vdots\\
a_{j1}&\cdots&a_{jn}\\
\vdots&&\vdots\\
a_{n1}&\cdots&a_{nn}
\end{matrix}\right|
+
\left|\begin{matrix}
a_{11}&\cdots&a_{1n}\\
\vdots&&\vdots\\
b_{j1}&\cdots&b_{jn}\\
\vdots&&\vdots\\
a_{n1}&\cdots&a_{nn}
\end{matrix}\right|.
\end{aligned}
\]
Với cột cũng có tính chất tương tự.

\paragraph{Hệ quả.}
Nhân một hàng (hoặc một cột) bất kỳ với số thực \(\lambda\), rồi cộng vào hàng
(hoặc cột) khác tương ứng, thì giá trị định thức không đổi. Ví dụ, khi
\(j\ne 1\),
\[
\left|\begin{matrix}
a_{11}+\lambda a_{j1}&\cdots&a_{1n}+\lambda a_{jn}\\
a_{21}&\cdots&a_{2n}\\
\vdots&&\vdots\\
a_{n1}&\cdots&a_{nn}
\end{matrix}\right|
=
\left|\begin{matrix}
a_{11}&\cdots&a_{1n}\\
a_{21}&\cdots&a_{2n}\\
\vdots&&\vdots\\
a_{n1}&\cdots&a_{nn}
\end{matrix}\right|.
\]

Dưới đây là một phương pháp tính định thức. Trước hết, không khó chứng minh
rằng với một ma trận tam giác,
\[
\left|\begin{matrix}
a_{11}&0&\cdots&0\\
a_{21}&a_{22}&\cdots&a_{2n}\\
\vdots&\vdots&\ddots&\vdots\\
a_{n1}&a_{n2}&\cdots&a_{nn}
\end{matrix}\right|
=
\left|\begin{matrix}
a_{11}&a_{12}&\cdots&a_{1n}\\
0&a_{22}&\cdots&a_{2n}\\
\vdots&\vdots&\ddots&\vdots\\
0&a_{n2}&\cdots&a_{nn}
\end{matrix}\right|
=
\prod_{i=1}^n a_{ii}.
\]

Xét định thức
\[
  \det B=
  \left|\begin{matrix}
  b_{11}&\cdots&b_{1n}\\
  \vdots&&\vdots\\
  b_{n1}&\cdots&b_{nn}
  \end{matrix}\right|.
\]
Ta có thể biến \(B\) thành dạng tam giác trên, ký hiệu \(B'\), đồng thời ghi
lại số lần đổi hàng là \(S\). Theo tính chất đổi hàng,
\[
  \det B=(-1)^S\det B'
  =
  (-1)^S
  \left|\begin{matrix}
  b'_{11}&b'_{12}&\cdots&b'_{1n}\\
  0&b'_{22}&\cdots&b'_{2n}\\
  \vdots&\ddots&\ddots&\vdots\\
  0&\cdots&0&b'_{nn}
  \end{matrix}\right|.
\]
Từ kết luận vừa nêu,
\[
  \det B=(-1)^S\prod_{i=1}^n b'_{ii}.
\]

Với ma trận \(A\), lấy \(2p\) số tự nhiên
\[
  1\le i_1<i_2<\cdots<i_p\le n,\qquad
  1\le j_1<j_2<\cdots<j_p\le n.
\]
Trong \(A\), lấy các phần tử giao giữa các hàng \(i_1,i_2,\ldots,i_p\) và các
cột \(j_1,j_2,\ldots,j_p\), giữ nguyên thứ tự ban đầu để tạo thành ma trận cấp
\(p\). Định thức của ma trận này gọi là một định thức con cấp \(p\) của \(A\),
ký hiệu
\[
  \minor{A}{i_1\ i_2\ \cdots\ i_p}{j_1\ j_2\ \cdots\ j_p}
  =
  \left|\begin{matrix}
  a_{i_1j_1}&\cdots&a_{i_1j_p}\\
  \vdots&&\vdots\\
  a_{i_pj_1}&\cdots&a_{i_pj_p}
  \end{matrix}\right|.
\]

Dựa trên định nghĩa này, ta giới thiệu một công thức rất quan trọng:
công thức Binet--Cauchy.

\paragraph{Định lý 8 (công thức Binet--Cauchy).}
Giả sử \(A\) là ma trận \(p\times q\) và \(B\) là ma trận \(q\times p\). Khi đó
\[
  \det(AB)=
  \begin{cases}
    0, & p>q,\\
    \det A\det B, & p=q,\\
    \displaystyle
    \sum_{1\le j_1<\cdots<j_p\le q}
    \minor{A}{1\ 2\ \cdots\ p}{j_1\ j_2\ \cdots\ j_p}
    \minor{B}{j_1\ j_2\ \cdots\ j_p}{1\ 2\ \cdots\ p},
    & p<q.
  \end{cases}
\]
Đặc biệt, với ma trận vuông \(A\),
\[
  \det(AA^{\mathsf T})=(\det A)^2.
\]

\section{Phương pháp mới}

\subsection{Giới thiệu}

Bây giờ ta giới thiệu một phương pháp mới: định lý Matrix-Tree, còn gọi là
định lý ma trận-cây Kirchhoff. Đây là một trong những công cụ mạnh nhất để xử
lý bài toán đếm cây khung. Định lý này được Kirchhoff chứng minh lần đầu vào
năm 1847. Trước khi phát biểu định lý, ta làm rõ vài khái niệm.

\begin{enumerate}
  \item Ma trận bậc của \(G\), ký hiệu \(D[G]\), là một ma trận \(n\times n\).
  Nó thỏa mãn: khi \(i\ne j\), \(d_{ij}=0\); khi \(i=j\), \(d_{ij}\) bằng bậc
  của đỉnh \(v_i\).

  \item Ma trận kề của \(G\), ký hiệu \(A[G]\), cũng là một ma trận
  \(n\times n\). Nó thỏa mãn: nếu giữa \(v_i\) và \(v_j\) có cạnh nối trực tiếp
  thì \(a_{ij}=1\), ngược lại \(a_{ij}=0\).
\end{enumerate}

Ta định nghĩa ma trận Kirchhoff của \(G\), còn gọi là toán tử Laplace, bởi
\[
  C[G]=D[G]-A[G].
\]
Định lý Matrix-Tree có thể phát biểu như sau: số cây khung khác nhau của
\(G\) bằng trị tuyệt đối của định thức của bất kỳ định thức con chính cấp
\(n-1\) nào của ma trận Kirchhoff \(C[G]\). Ở đây, định thức con chính cấp
\(n-1\) nghĩa là: với \(r\), \(1\le r\le n\), xóa đồng thời hàng thứ \(r\) và
cột thứ \(r\) của \(C[G]\), thu được ma trận mới ký hiệu \(C_r[G]\).

Ta xét một ví dụ để minh họa định lý Matrix-Tree. Giả sử \(G\) là đồ thị vô
hướng có \(5\) đỉnh như Hình a.

\begin{center}
  \small Hình a: đồ thị có các cạnh
  \((1,2),(1,3),(2,3),(2,4),(3,5),(4,5)\).
\end{center}

Theo định nghĩa, ma trận Kirchhoff \(C[G]\) của nó là
\[
  C[G]=
  \mat{
   2&-1&-1& 0& 0\\
  -1& 3&-1&-1& 0\\
  -1&-1& 3& 0&-1\\
   0&-1& 0& 2&-1\\
   0& 0&-1&-1& 2
  }.
\]
Lấy \(r=2\), ta có
\[
  C_2[G]=
  \mat{
   2&-1& 0& 0\\
  -1& 3& 0&-1\\
   0& 0& 2&-1\\
   0&-1&-1& 2
  }.
\]
Tính định thức:
\[
  \det C_2[G]
  =
  \left|\begin{matrix}
   2&-1& 0& 0\\
  -1& 3& 0&-1\\
   0& 0& 2&-1\\
   0&-1&-1& 2
  \end{matrix}\right|
  =11.
\]
Vì vậy đồ thị này có \(11\) cây khung.

\begin{center}
  \small Hình b: mười một cây khung khác nhau của đồ thị trong Hình a.
\end{center}

Định lý này trông rất ``kỳ diệu'': chỉ cần tính định thức của một ma trận là
có thể thu được số cây khung khác nhau của đồ thị. Vì sao lại như vậy? Ta thử
chứng minh định lý.

\subsection{Chứng minh}

Không khó nhận ra trọng tâm của định lý là ma trận Kirchhoff của đồ thị. Ma
trận này có tính chất đặc biệt gì? Qua phân tích, ta thấy:
\begin{enumerate}
  \item Với mọi đồ thị \(G\), định thức của ma trận Kirchhoff \(C\) luôn bằng
  \(0\), vì tổng các phần tử trên mỗi hàng và trên mỗi cột của \(C\) đều bằng
  \(0\).

  \item Nếu \(G\) không liên thông, thì định thức của mọi định thức con chính
  của ma trận Kirchhoff \(C\) đều bằng \(0\).
\end{enumerate}

Chứng minh tính chất thứ hai như sau. Nếu trong \(G\) có
\(k\), \(k>1\), thành phần liên thông \(G_1,G_2,\ldots,G_k\), ta có thể sắp xếp
lại các hàng và cột của \(C\) sao cho các đỉnh thuộc \(G_1\) xuất hiện trước,
rồi đến \(G_2\), v.v. Theo tính chất đổi hàng, giả sử ta thực hiện tổng cộng
\(t\) lần đổi hàng; hiển nhiên cũng cần \(t\) lần đổi cột. Tổng số lần đổi là
\(2t\), nên dấu của định thức không đổi.

Sau khi sắp xếp lại, ma trận có dạng khối chéo:
\[
  \mat{
  C[G_1]&0&\cdots&0\\
  0&C[G_2]&\cdots&0\\
  \vdots&\vdots&\ddots&\vdots\\
  0&0&\cdots&C[G_k]
  }.
\]
Các khối \(C[G_i]\) tương ứng với các thành phần liên thông. Ngoài các khối
này đều là \(0\), vì giữa hai thành phần liên thông khác nhau không có cạnh.

Nếu \(r\) là đỉnh thứ \(j\) trong thành phần liên thông thứ \(i\), thì
\[
  \det C_r[G]
  =
  \det C[G_1]\det C[G_2]\cdots
  \det C_j[G_i]\cdots
  \det C[G_k]
  =0.
\]

Tính chất tiếp theo là:

\begin{enumerate}[resume]
  \item Nếu \(G\) là một cây, thì định thức của mọi định thức con chính cấp
  \(n-1\) của ma trận Kirchhoff \(C\) đều bằng \(1\).
\end{enumerate}

Để chứng minh tính chất quan trọng này, ta liên tục biến đổi \(C_r[G]\) để thu
được một ma trận tam giác trên có toàn bộ phần tử trên đường chéo chính bằng
\(1\).

\begin{enumerate}
  \item Sắp xếp các đỉnh theo khoảng cách đến \(v_r\) trên cây, từ gần đến xa:
  trước hết là \(v_r\), sau đó là các đỉnh cách \(v_r\) một cạnh, tiếp theo là
  các đỉnh cách \(v_r\) hai cạnh, v.v. Các đỉnh có cùng khoảng cách có thể sắp
  tùy ý. Ta đánh số lại các đỉnh theo thứ tự này. Đồng thời, trong cây có gốc
  \(r\), gọi cha của \(v_i\) là \(v_j\). Rõ ràng trong thứ tự vừa thu được,
  \(v_j\) xuất hiện trước \(v_i\).

  \item Sắp xếp lại \(n-1\) hàng và \(n-1\) cột của \(C_r[G]\) theo thứ tự vừa
  thu được. Vì cần đổi cả hàng lẫn cột, tổng số lần đổi là số chẵn, nên định
  thức của \(C_r[G]\) không đổi.

  \item Xử lý theo thứ tự ngược lại của thứ tự trên. Với mỗi \(i\), nếu cha
  \(v_j\) của \(v_i\) không phải \(v_r\), cộng cột thứ \(i\) vào cột thứ
  \(j\).
\end{enumerate}

Sau khi xử lý như vậy, ta kiểm tra \(C_r[G]\) bằng quy nạp. Rõ ràng cột cuối
cùng thỏa yêu cầu, vì đỉnh cuối chỉ kề với cha của nó và có bậc \(1\). Giả sử
đang xử lý cột thứ \(i\), và các cột từ \(i+1\) đến \(n-1\) đều đã thỏa yêu
cầu. Gọi cha của \(v_i\) là \(v_j\), và các con của nó là
\[
  v_{i_1},v_{i_2},\ldots,v_{i_k}.
\]
Khi đó chỉ các cột \(i_1,i_2,\ldots,i_k\) mới được cộng vào cột \(i\). Mỗi
cột trong số đó có một phần tử \(-1\) ở hàng \(i\) và một phần tử \(1\) trên
đường chéo; còn cột \(i\) có giá trị \(-1\) ở các hàng
\(j,i_1,i_2,\ldots,i_k\), và có giá trị \(k+1\) ở hàng \(i\). Mỗi khi cộng một
cột con vào, phần tử \((i,i)\) giảm \(1\), đồng thời phần tử tương ứng ở hàng
con tăng lên thành \(0\). Cuối cùng, phần tử \((i,i)\) trở thành \(1\), còn
các hàng \(i_1,i_2,\ldots,i_k\) trong cột đó trở thành \(0\). Vì vậy
\(C_r[G]\) biến thành ma trận tam giác trên có toàn bộ đường chéo chính bằng
\(1\). Tính chất được chứng minh.

Trong chứng minh định lý Matrix-Tree, ta dùng ma trận liên thuộc \(B\) của đồ
thị vô hướng \(G\). Ma trận \(B\) có \(n\) hàng và \(m\) cột; hàng ứng với đỉnh,
cột ứng với cạnh. Nếu có cạnh \(e=\{v_i,v_j\}\), thì trong cột ứng với \(e\),
hai hàng tương ứng với \(v_i\) và \(v_j\) nhận một giá trị \(1\) và một giá trị
\(-1\), còn các hàng khác đều là \(0\). Việc chọn đỉnh nào nhận \(1\) và đỉnh
nào nhận \(-1\) không quan trọng.

\begin{center}
  \small Hình: ma trận liên thuộc của đồ thị vô hướng; mỗi cột có đúng một
  \(1\) và một \(-1\) tại hai đầu mút của cạnh.
\end{center}

Tiếp theo xét \(BB^{\mathsf T}\). Dễ chứng minh
\((BB^{\mathsf T})_{ij}\) bằng tích vô hướng của hàng thứ \(i\) và hàng thứ
\(j\) của \(B\). Vì vậy, khi \(i=j\), \((BB^{\mathsf T})_{ij}\) bằng số cạnh
nối với \(v_i\), tức bậc của \(v_i\). Khi \(i\ne j\), nếu có một cạnh nối
\(v_i\) và \(v_j\), thì \((BB^{\mathsf T})_{ij}=-1\), ngược lại bằng \(0\).
Điều này hoàn toàn trùng với định nghĩa của ma trận Kirchhoff. Do đó
\[
  C=BB^{\mathsf T}.
\]
Nói cách khác, ta có thể chuyển vấn đề về \(C\) sang \(BB^{\mathsf T}\). Việc
này có ích gì?

Gọi \(B_r\) là ma trận nhận được từ \(B\) sau khi xóa hàng thứ \(r\). Khi đó
\[
  C_r=B_rB_r^{\mathsf T}.
\]
Theo công thức Binet--Cauchy,
\[
\begin{aligned}
  \det(C_r)
  &=\det(B_rB_r^{\mathsf T})\\
  &=
  \sum_{\substack{x\subseteq E\\ |x|=n-1}}
  \det(B_r^x)\det\!\bigl((B_r^x)^{\mathsf T}\bigr)\\
  &=
  \sum_{\substack{x\subseteq E\\ |x|=n-1}}
  \det\!\bigl(B_r^x(B_r^x)^{\mathsf T}\bigr)
  =
  \sum_{\substack{x\subseteq E\\ |x|=n-1}}
  \bigl(\det B_r^x\bigr)^2.
\end{aligned}
\]
Ở đây \(B_r^x\) là ma trận tạo bởi các cột của \(B_r\) tương ứng với các cạnh
thuộc \(x\).

Ta chú ý rằng nếu các cạnh trong \(x\) tạo thành một chu trình thì luôn có
\[
  \det B_r^x=0.
\]
Chẳng hạn, nếu đồ thị mới có một chu trình độ dài \(3\), ta có thể sắp xếp lại
\(B_r^x\) sao cho khối liên quan đến chu trình nằm ở góc trái trên.

\begin{center}
  \small Hình: sau khi sắp xếp lại \(B_r^x\), các hàng và cột của một chu trình
  tạo thành một khối suy biến ở góc trái trên.
\end{center}

Khi đó \(\det B_r^x\) bằng định thức của định thức con cấp \(3\) ở góc trái
trên nhân với định thức của ma trận cấp \(n-4\) ở góc phải dưới. Định thức con
cấp \(3\) này suy biến, vì tổng trên mỗi hàng và mỗi cột của nó đều bằng
\(0\). Do đó \(\det B_r^x=0\). Lập luận tương tự có thể mở rộng cho chu trình
có độ dài bất kỳ.

Rõ ràng \(B_r^x(B_r^x)^{\mathsf T}\) có thể xem là một định thức con chính cấp
\(n-1\) của ma trận Kirchhoff của đồ thị mới chỉ gồm toàn bộ các đỉnh và các
cạnh thuộc \(x\). Theo các tính chất của ma trận Kirchhoff, nếu thêm tất cả
\(n-1\) cạnh thuộc \(x\) vào đồ thị tạo thành một cây, thì
\[
  \det\!\bigl(B_r^x(B_r^x)^{\mathsf T}\bigr)=1;
\]
nếu không tạo thành cây, thì chắc chắn tồn tại chu trình, và
\(\det B_r^x=0\). Nói cách khác, ta xét tất cả các tập con cạnh có kích thước
\(n-1\): nếu tập con đó tạo thành một cây thì đáp án tăng thêm \(1\), còn nếu
không thì không đóng góp. Điều này đúng bằng số cây khung của đồ thị ban đầu.
Vì vậy ta đã chứng minh được định lý Matrix-Tree.

Nếu đồ thị có cạnh song song, định lý Matrix-Tree vẫn đúng; cách chứng minh
tương tự và dành cho người đọc tự suy nghĩ. Vì tính định thức mất
\(\Oh(n^3)\) thời gian, bài toán đếm cây khung cũng có thể hoàn thành trong
\(\Oh(n^3)\) thời gian.

\subsection{Cách hiểu trực quan}

Phân tích vừa rồi có thể hơi sâu. Dưới đây ta nhìn định lý này theo cách trực
quan hơn.

Trước hết xét một bài toán toán học đơn giản: hãy tìm số nghiệm nguyên không
âm của phương trình
\[
  x_1+x_2+x_3=2.
\]
Đây là một bài toán đếm tổ hợp quen thuộc. Cách làm thường dùng là lấy hai số
\(1\) và hai dấu phân cách \(\triangle\), rồi sắp xếp tùy ý bốn phần tử này.
Số cách sắp xếp khác nhau đúng bằng số nghiệm của phương trình, tức
\[
  \binom{4}{2}.
\]
Vì sao làm như vậy?

Liệt kê \(6\) cách sắp xếp, ta thấy mỗi cách tương ứng với một nghiệm của
phương trình:
\[
\begin{array}{rcl}
  \triangle\triangle 11 &\leftrightarrow& x_1=0,\ x_2=0,\ x_3=2,\\
  \triangle 1\triangle 1 &\leftrightarrow& x_1=0,\ x_2=1,\ x_3=1,\\
  \triangle 11\triangle &\leftrightarrow& x_1=0,\ x_2=2,\ x_3=0,\\
  \multicolumn{3}{c}{\cdots}
\end{array}
\]

Nói cách khác, thông qua chuyển đổi mô hình, ta tìm ra quan hệ tương ứng giữa
bài toán ban đầu và bài toán mới, rồi dùng kiến thức toán học liên quan để
giải bài toán mới; nhờ đó giải luôn bài toán ban đầu. Ý nghĩa quan trọng của
sự chuyển đổi này là nó dựng nên cầu nối giữa các bài toán khác nhau.

Quay lại định lý Matrix-Tree. Ta cũng chuyển đổi mô hình, từ mô hình đồ thị
sang mô hình ma trận, rồi phát hiện mỗi hạng tử trong khai triển theo công
thức Binet--Cauchy tương ứng với một tập con cạnh có kích thước \(n-1\). Trong
đó, hạng tử có giá trị \(1\) tương ứng với một cây khung; các hạng tử không
tương ứng với cây khung có giá trị \(0\). Như vậy, bài toán được chuyển thành
tính tổng tất cả các hạng tử trong khai triển. Dùng kiến thức toán học sẵn có
là có thể giải quyết thành công.

So sánh hai bài toán, ta thấy: trong bài toán thứ nhất, nghiệm của phương
trình đóng vai trò giống cây khung của đồ thị trong bài toán thứ hai; phương
án sắp xếp trong bài toán thứ nhất đóng vai trò giống mỗi hạng tử trong khai
triển ở bài toán thứ hai. Đồng thời, trong mỗi bài toán, hai đối tượng lại có
quan hệ tương ứng với nhau. Khác biệt chỉ nằm ở chỗ quan hệ tương ứng trong
bài toán thứ hai kín đáo và phức tạp hơn, cần một lý thuyết toán học mạnh hơn
để chống đỡ.

\section{Ứng dụng cụ thể}

Dưới đây ta xem một vài ví dụ về bài toán đếm cây khung xuất hiện trong thi
tin học.

\subsection*{Ví dụ 2: Organising the Organisation (UVA p10766)}

Jimmy phụ trách việc phân cấp nhân sự trong công ty. Gần đây anh gặp một rắc
rối nhỏ: để nâng cao hiệu quả làm việc, hội đồng quản trị quyết định phân cấp
lại toàn bộ nhân viên.

\begin{center}
  \small Hình a: sơ đồ phân cấp nhân viên; \(A\) trực tiếp lãnh đạo \(B\) và
  \(D\), còn \(D\) trực tiếp lãnh đạo \(C\).
\end{center}

Tình hình sau khi phân cấp được mô tả như Hình a. Cụ thể, ngoại trừ một tổng
giám đốc, mỗi nhân viên khác có đúng một cấp trên trực tiếp. Do quan hệ cá
nhân giữa các nhân viên, có thể xuất hiện trường hợp \(a\) và \(b\) đều không
muốn người kia trở thành cấp trên trực tiếp của mình. Công ty có \(n\) nhân
viên, đánh số từ \(1\) đến \(n\), và hội đồng quản trị đã quyết định để nhân
viên số \(k\) làm tổng giám đốc. Nhiệm vụ của Jimmy là tính có bao nhiêu sơ đồ
phân cấp nhân viên khác nhau.

Quy mô dữ liệu: \(1\le n\le 50\).

\subsubsection*{Phân tích}

Cách giải bài này khá rõ. Nếu giữa \(a\) và \(b\) không có mâu thuẫn trực tiếp,
ta nối một cạnh giữa họ. Hiển nhiên đồ thị quan hệ nhân viên cuối cùng phải là
một cây khung của đồ thị ban đầu. Dù đề bài đã quy định gốc của cây khung,
số cây khung của đồ thị vô hướng không phụ thuộc vào gốc, nên ta chỉ cần trực
tiếp dùng định lý Matrix-Tree để tính số cây khung của đồ thị ban đầu.

\subsection*{Ví dụ 3: Nỗi phiền của nhà vua}

Quốc vương Byteotia của Byteland gần đây rất phiền lòng: đất nước của ông bị
lũ lụt tấn công. Trận lũ trăm năm mới gặp đã phá hủy nhiều con đường quan
trọng của Byteland, khiến cả quốc gia rơi vào tê liệt. Byteland gồm \(n\)
thành phố, và giữa hai thành phố khác nhau bất kỳ đều có đường nối. Để khôi
phục trật tự bình thường càng nhanh càng tốt, Byteotia tổ chức cho các chuyên
gia nghiên cứu và liệt kê tất cả các con đường có thể đi lại bình thường
(những con đường khác có thể vẫn đang ngập trong nước). Trong số đó, có đường
đã bị phá hủy và cần sửa lại; có đường vẫn có thể tiếp tục sử dụng. Điều kỳ
lạ là tất cả các con đường còn có thể sử dụng không tạo thành chu trình. Để
tiết kiệm thời gian tối đa, Byteotia muốn chỉ sửa số đường ít nhất sao cho cả
quốc gia liên thông. Ban đầu Byteotia định đánh giá từng phương án rồi chọn
phương án tối ưu, nhưng rất nhanh ông phát hiện số phương án quá lớn. Vì vậy,
ông tìm đến bạn, chuyên gia máy tính giỏi nhất của Byteland, nhờ bạn tính số
phương án sửa đường có thể có.

Quy mô dữ liệu: \(1\le n\le 500\).

\subsubsection*{Phân tích}

Thoạt nhìn bài này khá khó, vì phải đếm số phương án sửa ít đường nhất, đồng
thời còn phải xét các con đường không cần sửa. Phân tích kỹ hơn, ta thấy một
điều kiện rất quan trọng: tất cả các con đường còn có thể sử dụng không tạo
thành chu trình. Điều này tạo thuận lợi cho lời giải. Ta biết rằng để một đồ
thị \(n\) đỉnh liên thông, ít nhất cần \(n-1\) cạnh; các cạnh đó tạo thành một
cây. Một tính chất quan trọng của cây là không có chu trình, nên mọi con đường
vẫn dùng được chắc chắn phải xuất hiện trong cây cuối cùng. Như vậy, ta đã
chuyển bài toán thành: đếm số cây khung của một đồ thị, trong đó một số cạnh
bắt buộc phải chọn.

Đây cũng là một bài toán đếm cây khung. Tuy nhiên, do tồn tại các cạnh bắt
buộc, ta không thể áp dụng trực tiếp định lý Matrix-Tree. Phải làm gì?

Ta biết rằng nếu yêu cầu cây khung phải chứa một cạnh \(e\), thì số cây khung
\(t(G)\) của đồ thị \(G\) bằng \(t(G/e)\), tức số cây khung của đồ thị mới
nhận được sau khi co cạnh \(e\). Vì vậy cần:
\begin{enumerate}
  \item co tất cả các cạnh bắt buộc;
  \item tính số cây khung của đồ thị mới sau khi co.
\end{enumerate}

Co một cạnh mất \(\Oh(n)\) thời gian, và nhiều nhất chỉ có \(n-1\) cạnh cần
co, nên bước này có độ phức tạp \(\Oh(n^2)\). Theo phân tích phía trên, tính
số cây khung của một đồ thị có độ phức tạp \(\Oh(n^3)\). Vì vậy ta giải được
toàn bộ bài toán trong thời gian \(\Oh(n^3)\).

\section{Tổng kết}

Bài viết bắt đầu từ một ví dụ trong thi tin học, giới thiệu thuật toán đếm cây
khung, tư tưởng toán học bên trong thuật toán, và ứng dụng của nó trong thi
tin học. Trong quá trình thảo luận, ta liên hệ chặt chẽ hai đối tượng tưởng
như không liên quan: cây khung của đồ thị và định thức của ma trận. Trước hết,
ta chuyển bài toán đếm cây khung của đồ thị thành bài toán tính định thức ma
trận; sau đó dùng chứng minh toán học để luận chứng quan hệ tương ứng giữa
định thức ma trận và cây khung của đồ thị. Từ đó không khó nhận ra: nền tảng
toán học vững chắc là bảo đảm để giải quyết vấn đề, còn liên tưởng sáng tạo
là linh hồn của lời giải.

\section*{Tài liệu tham khảo}

\begin{enumerate}
  \item Ju Yuma và các tác giả khác, \textit{Linear Algebra}.
  \item Xu Yichao, \textit{Introduction to Algebra}.
  \item Reinhard Diestel, \textit{Graph Theory II}.
  \item Hu Maolin, \textit{Counting spanning trees in connected graphs
  containing specified edges}.
  \item Martin Rubey, \textit{Counting Spanning Tree}.
\end{enumerate}

\end{document}
